\documentclass[11pt, addpoints]{modelo/prova2}
\usepackage[T1]{fontenc}
\usepackage[brazilian]{babel}

\include{orientacoes}
\begin{document}

\maketitle 
\question[1.5] Sobre fundamentos de redes e meios físicos, responda ao que se pede:
\begin{choices}
    \item Explique a diferença entre os conceitos de largura de banda (bandwidth) e taxa de transmissão (throughput) em uma rede de computadores. 
    \item Explique os quatro tipos de atraso no envio de um pacote em uma rede de computadores.
    \item Em um ambiente hospitalar com equipamentos eletrônicos que geram campos magnéticos intensos, qual meio físico (par trançado ou fibra óptica) é mais indicado? Justifique com base nas propriedades físicas do meio. 
\end{choices}
\question[2] Uma empresa de 20 funcionários utiliza um Hub para interconectar todos os computadores e enfrenta lentidão com frequência. Foi proposta a troca por um Switch de Camada 2.
\begin{choices}
    \item Explique a diferença fundamental entre Hub e Switch no encaminhamento de quadros e como isso impacta os domínios de colisão.
    \item Descreva como o Switch constrói dinamicamente sua tabela de endereços MAC e o que ocorre ao receber um quadro cujo endereço MAC de destino não está na tabela. 
    \item O que são VLANs e como podem contribuir para a segurança e o desempenho dessa empresa?
\end{choices}

\question[1.5] A respeito de redes sem fio (IEEE 802.11), discorra sobre os itens a seguir:
\begin{choices}
    \item Diferencie o funcionamento do CSMA/CD (Ethernet) e do CSMA/CA (Wi-Fi). Por que as redes sem fio não utilizam detecção de colisão durante a transmissão?
    \item Explique o problema do terminal escondido e como o mecanismo RTS/CTS pode mitigá-lo.
    \item Por que aplicações de tempo real como videoconferência tendem a sofrer mais em redes Wi-Fi do que em redes Ethernet cabeadas? Relacione sua resposta ao mecanismo de retransmissão na camada de enlace. 
\end{choices}
\question[1.5] No que se refere ao Modelo OSI, Arquitetura TCP/IP e VLANs, resolva:
\begin{choices}
    \item Compare os Modelos OSI e TCP/IP quanto ao número de camadas e às funções das camadas de Transporte e Rede (Internet).
    \item Explique por que o protocolo IP é classificado como ``melhor esforço'' (best-effort) e como o protocolo TCP atua para garantir a confiabilidade na entrega dos dados fim-a-fim. 
    \item Qual a função do padrão IEEE 802.1Q e como ele permite que múltiplas VLANs trafeguem por um mesmo enlace físico entre switches?
\end{choices}
\newpage
\question[2.5] Um administrador de rede precisa dividir o bloco 192.168.30.0/24 para atender a três departamentos:
\begin{itemize}
    \item Marketing: 55 hosts
    \item Desenvolvimento: 28 hosts
    \item Suporte: 8 hosts
\end{itemize}
Utilizando a técnica de VLSM (Variable Length Subnet Mask), determine:
\begin{choices}
    \item A máscara de sub-rede (em notação decimal) e o prefixo (em notação /X) para cada departamento, justificando a escolha com base na quantidade de bits necessários para a parte de host.
    \item O intervalo de endereços IP utilizáveis (primeiro e último host válido) e o endereço de broadcast de cada sub-rede. 
    \item A quantidade total de endereços desperdiçados (não utilizados) após a divisão.
\end{choices}
\question[2.0] Acerca de Roteamento: IGP, EGP e Planos da Camada de Rede, elabore sua resposta sobre:
\begin{choices}
    \item Diferencie protocolos IGP (como OSPF e RIP) de protocolos EGP (como BGP) quanto aos seus objetivos, métricas e alcance.
    \item Explique a diferença entre Roteamento (Routing) e Encaminhamento (Forwarding), indicando em qual plano — Plano de Controle ou Plano de Dados — cada função atua.
    \item Em cenários reais da Internet, o BGP frequentemente prioriza políticas e acordos comerciais em vez de escolher o caminho tecnicamente mais curto. Explique por que isso ocorre e dê um exemplo de situação em que uma rota pode ser rejeitada mesmo sendo a mais curta. 
\end{choices}

\end{document}